\section*{Chapitre \ref{ch:mx:market-quality-and-regulation} --- Qualité de marché et régulation}

\begin{solution}{exo:mx:market-quality-and-regulation:1}
$(3.18-1.58)-(2.27-1.56)=1.60-0.71=0.89$ points de base.
\end{solution}

\begin{solution}{exo:mx:market-quality-and-regulation:2}
$900/250=3.6$ messages par transaction~; avec l'effet de la règle, environ 2.4.
\end{solution}

\begin{solution}{exo:mx:market-quality-and-regulation:3}
Les transactions alternent entre cours acheteur et cours vendeur, si bien que les rendements sur de courts intervalles sont autocorrélés négativement et que leurs variances s'additionnent pour dépasser celle de l'intervalle long.
La règle a abaissé le ratio de 0.14 (0.13 estimé)~: fourchettes plus larges, rebond bid-ask plus fort.
\end{solution}

\begin{solution}{exo:mx:market-quality-and-regulation:4}
Les journées d'une même action partagent ses caractéristiques non observables et leurs erreurs sont corrélées dans le temps~; la règle est assignée par action, si bien que le nombre effectif
d'observations indépendantes est le nombre d'actions, et non celui des couples action-jour.
\end{solution}

\begin{solution}{exo:mx:market-quality-and-regulation:5}
Il visait 9~\% du volume de la minute précédente, et le \omterm{def:mx:market-quality-and-regulation:hotpotato}{trading de la patate chaude} entre intermédiaires a gonflé ce volume sans en absorber la moindre part des ventes~:
l'algorithme vendait plus vite sur un marché qui n'achetait pas.
\end{solution}

\begin{solution}{exo:mx:market-quality-and-regulation:6}
Les données du carnet d'ordres ne montrent que les ordres qui ont modifié le carnet~: les perdants d'une course (tentatives échouées et annulations) ne laissent aucune trace. Les données de messages de la bourse,
avec chaque message et son issue, en gardent la trace.
\end{solution}

\begin{solution}{exo:mx:market-quality-and-regulation:7}
1.11 points de base (erreur type 0.06), contre une valeur vraie de 0.94~: les actions les moins actives subissent davantage le changement de marché que les actions du groupe témoin,
si bien que le groupe témoin sous-estime ce qu'auraient fait les fourchettes des actions traitées sans la règle~; la condition de tendances parallèles ne tient pas.
\end{solution}

\begin{solution}{exo:mx:market-quality-and-regulation:8}
Les 1.6 points de base incluent le changement de marché qui a touché toutes les actions (0.71 sur les témoins)~; l'effet propre de la règle, par \omterm{def:mx:tick-size-queues-and-priority:did}{doubles différences},
vaut 0.89, proche de la valeur vraie de 0.87.
\end{solution}

\begin{solution}{pb:mx:market-quality-and-regulation:1}
\textbf{1.} Coûts, profondeur et contenu informatif des prix~: \omterm{def:mx:decomposing-the-spread:quoted}{fourchettes affichées}, effectives et réalisées, impact, profondeur, résilience, Amihud, ratio de variance.
\textbf{2.} Fourchette effective = fourchette réalisée + impact sur le prix.
\textbf{3.} Voir les définitions.
\textbf{4.} Écart des prix à une marche aléatoire sur de courts intervalles~: inférieur à un en présence de rebond bid-ask.
\textbf{5.} Le \omterm{def:mx:market-quality-and-regulation:hft}{trading algorithmique} a resserré les fourchettes, réduit la sélection adverse et rendu les cotations plus informatives pour les grandes actions~; effet identifié grâce à
l'automatisation de la diffusion des cotations du NYSE en 2003, utilisée comme instrument.
\textbf{6.} Les traders à haute fréquence négocient dans le sens des variations permanentes des prix et contre les erreurs transitoires, surtout par leurs ordres agressifs.
\textbf{7.} 75\,000 contrats E-mini à 9~\% du volume, sans égard au prix ni au temps~; les traders à haute fréquence ont ensuite échangé plus de 27\,000 contrats entre
eux en 14 secondes, soit environ 49~\% du volume, en achetant environ 200 contrats nets.
\textbf{8.} Courses d'arbitrage de latence (environ une par minute et par titre du FTSE~100, durée modale de 5--10 microsecondes, environ 20~\% du volume), observées dans les données de messages
qui enregistrent les tentatives échouées.
\textbf{9.} 20 actions, 16 journées de dix minutes~; à partir du jour 8, une hausse de la volatilité et une baisse des cotations pour toutes, ainsi qu'un plafond du \omterm{def:mx:market-quality-and-regulation:otr}{ratio ordres sur transactions} pour
chaque autre action.
\textbf{10.} En simulant chaque couple action-jour traité après l'événement, également sans la règle, avec le même flux d'ordres.
\textbf{11.} $y_{id}=\alpha_i+\gamma_d+\beta\,\text{traité}_i\times\text{après}_d+\varepsilon_{id}$, erreurs regroupées par action.
\textbf{12.} \emph{Résultat principal}~: \omterm{def:mx:tick-size-queues-and-priority:did}{doubles différences} de 0.89 point de base (erreur type 0.09), contre une comparaison avant-après de 1.60 (0.08) qui ignore
le groupe témoin~; la valeur vraie est 0.87 (0.03).
\textbf{13.} Coefficients à moins de 0.07 de zéro avant l'événement~: les groupes évoluaient ensemble, ce qui soutient l'hypothèse de tendances parallèles.
\textbf{14.} Elle a réduit les messages mais élargi les fourchettes (0.9 affichée, 1.3 effective), réduit la profondeur et accru le rebond bid-ask~: pire selon chaque mesure du
coût de négociation.
\textbf{15.} Les \omterm{def:mx:tick-size-queues-and-priority:did}{doubles différences} surestiment~: 1.11 contre 0.94.
\textbf{16.} Le \omterm{def:mx:market-quality-and-regulation:otr}{ratio ordres sur transactions} (en baisse d'environ un quart), cible de la règle.
\textbf{17.} Un groupe témoin comparable (ou une adoption échelonnée), des données de la période précédente pour vérifier les tendances, et des mesures issues des données au niveau des messages.
\textbf{18.} Le \omterm{def:mx:market-quality-and-regulation:hotpotato}{trading de la patate chaude} le gonfle sans que personne ne prenne de risque, comme lors du flash crash.
\textbf{19.} Les plateformes de négociation doivent calculer, pour chaque instrument, le ratio des ordres non exécutés sur les transactions de chaque membre.
\textbf{20.} Une mesure de ce qui serait arrivé de toute façon.
\end{solution}

\begin{solution}{iq:mx:market-quality-and-regulation:1}
Comparer avec des actions que le changement n'a pas touchées sur les mêmes journées, vérifier qu'elles évoluaient ensemble auparavant, et contrôler la volatilité et le volume~; une méthode de
\omterm{def:mx:tick-size-queues-and-priority:did}{doubles différences} couplée à une étude d'événement.
\iqlookfor{Un groupe témoin~; tendances préalables.}
\end{solution}

\begin{solution}{iq:mx:market-quality-and-regulation:2}
La variation du groupe traité moins celle du groupe témoin~; elle suppose que les deux auraient évolué de la même façon sans le traitement (tendances parallèles), testée par
les coefficients de la période précédente de l'étude d'événement et par des dates placebo.
\iqlookfor{Tendances parallèles~; tests placebo.}
\end{solution}

\begin{solution}{iq:mx:market-quality-and-regulation:3}
En temps normal, plutôt oui~: fourchettes plus étroites, prix plus informatifs (Hendershott, Jones et Menkveld~; Brogaard, Hendershott et Riordan). Il ne
règle pas le comportement en période de tension (le flash crash), le coût des courses de latence (Aquilina, Budish et O'Neill), ni qui supporte la sélection adverse.
\iqlookfor{Bilan~; tensions et courses.}
\end{solution}

\begin{solution}{iq:mx:market-quality-and-regulation:4}
Un \omterm{def:mx:benchmark-algorithms:vwap}{algorithme de participation} ou de suivi du volume a besoin d'une limite de variation de cours et d'un contrôle que le volume correspond à une vraie liquidité~; surveiller la profondeur du carnet d'ordres, et pas seulement le volume,
et ralentir lorsque les prix évoluent vite.
\iqlookfor{Sensibilité au prix~; volume contre profondeur.}
\end{solution}

\begin{solution}{iq:mx:market-quality-and-regulation:5}
Compter les messages et les transactions par instrument et par compte dans la passerelle d'ordres, calculer le ratio de la plateforme selon sa propre formule en cours de séance, brider le débit de
stratégies de cotation à mesure qu'elles approchent de la limite, et alerter avant les dépassements.
\iqlookfor{Comptage en temps réel~; limitation du débit~; la définition de la plateforme.}
\end{solution}

\begin{solution}{iq:mx:market-quality-and-regulation:6}
À partir des données de messages, avec chaque ordre, chaque annulation et chaque tentative échouée, et des horodatages de la bourse~: regrouper les messages qui touchent le même \omterm{def:mx:the-limit-order-book:tick}{niveau de prix} dans
quelques microsecondes, avec au moins un gagnant et un perdant~; la plateforme doit fournir les messages échoués, et pas seulement le carnet d'ordres.
\iqlookfor{Données au niveau des messages~; définir une course.}
\end{solution}
